\documentclass[11pt]{article}

\usepackage[a4paper, margin=1in]{geometry}
\usepackage{parskip}
\usepackage{hyperref}
\usepackage{microtype}
\usepackage{titlesec}

\titleformat{\section}{\large\bfseries}{}{0pt}{}
\titleformat{\subsection}{\normalsize\bfseries}{}{0pt}{}

\title{What we did, in plain English: \\
\large Using modern AI to understand endometriosis patients' Reddit posts}
\author{[Author] (mentored by [Mentor])}
\date{\today}

\begin{document}
\maketitle

\section{The big picture}

Endometriosis is a long-lasting condition where tissue similar to the
lining of the uterus grows in places it shouldn't, like the ovaries or
the abdomen. It causes severe pain, can hurt fertility, and on average
takes about seven to ten years to even get diagnosed~\cite{young2015}.
Doctors often don't have a clear answer for what works, and patients
end up sharing what they go through with each other on the
internet---especially on Reddit, where there are entire communities
dedicated to endometriosis~\cite{eysenbach2009infodemiology}.

A 2024 research abstract showed something interesting: you can take
thousands of Reddit posts written by patients and feed them into a
large language model (LLM) like GPT-4, and the model will summarize
what people are saying about their care. That earlier study found five
big themes about endometriosis treatment.

Our project asks a narrower question: can a completed
single-extractor LLM workflow produce a larger, auditable theme list
from an r/Endo archive, and can we show where that theme list is
well-supported or weak?

The completed run is not a two-extractor replication. It is a
transparent single-extractor run with automated checks for quote
grounding, literature alignment, and judge-model ratings.

\section{Where the data came from}

We used a mentor-provided archive from \texttt{r/Endo} only. The archive
contains 38{,}001 posts and 291{,}754 comments. The completed run does
not use a public-archive retrieval and does not make other-subreddit,
year, or subgroup claims.

Once we have the raw data, we save it as JSON files on disk and
compute a SHA-256 ``fingerprint'' of every file. That fingerprint
guarantees that if anyone re-runs our analysis, they're using exactly
the same data we did---if even one character were different, the
fingerprint would change. This is how we make the science
reproducible.

For the actual analysis we use a sample of 2{,}000 posts and
80{,}000 comments from that archive.

\subsection{Ethics}

Reddit posts are public, but the people writing them aren't expecting
to show up in a research paper word-for-word. So in our published
manuscript, every quote is paraphrased, not copy-pasted, and we don't
include any usernames or direct messages~\cite{norori2021addressing}.
This kind of work is normally exempt from IRB review (institutional
ethics review) because it uses public data and doesn't identify
individuals.

\section{The AI workflow}

The 2024 abstract used GPT-4. We cannot re-run that exact analysis, so
we treat its five themes as a frozen baseline. Our completed run uses
Claude Haiku 4.5 as the only theme extractor~\cite{anthropic_claude_haiku}.
Claude Sonnet~4.6 reduces the chunk-level outputs into final themes,
and Claude Sonnet~4.6 plus Claude Opus~4.7 judge those themes with a
rubric.

\subsection{How we feed posts to the AI}

A whole archive of Reddit posts is way too long to fit into an AI's
memory at once. So we use a pattern called ``map-reduce'':

\begin{enumerate}
  \item Split the sample into 74 chunks of about 100{,}000 tokens each,
    with 2{,}000 tokens of overlap between neighboring chunks.
  \item Ask Claude Haiku 4.5, using prompt~A at temperature $T=0.0$,
    to summarize each chunk. This is the ``map'' step, and it produced
    722 chunk-level themes.
  \item Ask a Claude Sonnet~4.6 subagent to merge those chunk-level
    themes into one global list. This is the ``reduce'' step, and it
    produced 20 final themes.
\end{enumerate}

Prompt~B, Gemini extraction, and higher-temperature stability runs were
not completed, so we do not report them as evidence.

\section{How we checked the AI's answers}

Instead of just trusting the AI, we score each final theme using the
checks that actually completed.

\paragraph{1. Grounding (faithfulness).} Every time the AI gives us a
theme, it cites a few quotes from real posts. We check whether each
quote appears in the corpus after simple text normalization. If not, a
BM25 keyword search looks for the closest source passage. The fraction
of the AI's quotes that pass this check is the ``grounding rate.'' A
low grounding rate suggests the AI is hallucinating.

\paragraph{2. Anchored validity (Young 2015).} There's already a
respected systematic review of endometriosis patient experiences:
Young, Fisher, and Kirkman (2015), which carefully reviewed dozens of
qualitative studies~\cite{young2015}. We encoded that paper's main
themes (diagnostic delay, chronic pain, treatment disappointment,
fertility worries, etc.) as a JSON anchor. If our AI finds a theme
that matches Young 2015, that's evidence the AI is in line with the
peer-reviewed literature. If it finds a theme Young didn't catch,
that's a candidate ``new finding'' worth a closer look.

\paragraph{3. AI-as-judge with a rubric.} We let Claude Sonnet~4.6 and
Claude Opus~4.7 score the final themes on five dimensions, each from 1 to 5:
faithfulness, comprehensiveness, clinical relevance, specificity, and
non-redundancy. These models are judges, not second extractors.

\paragraph{4. Chunk recurrence.} We also record how many of the 74 map
chunks support each final theme. This helps show whether a theme came
from many parts of the sample, but it is not the same as cross-model
agreement.

\subsection{Combining everything into one number}

The completed composite reliability score averages the available
normalized signals: grounding, Young 2015 anchor alignment, and
Sonnet/Opus judge scores. A theme with a higher score is one that:

\begin{itemize}
  \item Has quotes that actually exist in the corpus
  \item Matches a known theme from prior research
  \item Got high judge scores from Sonnet and Opus
\end{itemize}

A theme near 0.0 is one that probably came from one AI's quirks rather
than from the corpus.

We follow the MI-CLAIM checklist for clinical AI
papers~\cite{norgeot2020miclaim} so reviewers can see we've thought
through the standard reporting items.

\section{What the completed run produced}

The completed run is a single-extractor/subagent result: Haiku for the
map step, Sonnet for the reduce step, and Sonnet plus Opus for judging.

What we ended up with:

\begin{itemize}
  \item \textbf{Corpus}: 38{,}001 posts and 291{,}754 comments in the
    mentor-provided \texttt{r/Endo} archive; 2{,}000 posts and
    80{,}000 comments in the analysis sample.
  \item \textbf{Map step}: Claude Haiku 4.5, prompt~A, $T=0.0$, 74
    subagents, 100{,}000-token chunks, 2{,}000-token overlap, and 722
    chunk-level themes.
  \item \textbf{Reduce step}: Claude Sonnet~4.6 merged those outputs
    into \textbf{20 final themes}.
  \item \textbf{Quote grounding}: 21 of 60 supporting quotes were
    grounded, or \textbf{35\%}.
  \item \textbf{Young 2015 alignment}: 9 of 20 final themes matched a
    Young 2015 anchor, covering 7 of the 8 anchors.
  \item \textbf{Judges}: Claude Sonnet~4.6 and Claude Opus~4.7 scored
    the final themes with the rubric.
  \item \textbf{Composite reliability}: mean \textbf{0.55}, range
    0.28--0.87.
\end{itemize}

\section{What this means}

Here's why this work matters, in plain terms.

First, it shows what a single-extractor LLM can and cannot support.
The 20 final themes are not just free-form AI prose: each one carries
grounding, anchor, judge, and chunk-recurrence evidence.

Second, we now have a way to check whether AI summaries are honest. A
``grounding rate'' tells you how often the AI's quotes actually exist
in the source. This is the kind of basic check that most LLM papers
in healthcare don't do, and reviewers should start expecting it.

Third, by anchoring our findings to the Young 2015 systematic review,
we can separate ``things AI found that scientists already
knew''---which is reassuring---from ``things AI found that scientists
hadn't catalogued yet''---which is where the next research questions
live.

Fourth, the pipeline is reproducible in the limited sense that the data
sample, prompts, model snapshots, and outputs are recorded. It is not a
claim of cross-extractor reproducibility.

\section{What we couldn't do}

We're upfront about the limits:

\begin{itemize}
  \item \textbf{Reddit isn't representative.} People who post on
    \texttt{r/Endo} tend to be young, English-speaking, and online.
    Their experiences don't speak for everyone with the condition.
  \item \textbf{There is only one extractor.} Gemini extraction did not
    complete, so this paper does not report cross-model agreement.
  \item \textbf{Prompt-B and high-temperature stability are absent.}
    The completed run is prompt~A at $T=0.0$ only.
  \item \textbf{We don't have a working OpenAI key}, so we can't
    re-run the 2024 GPT-4 baseline ourselves. The 2024 abstract's
    themes are treated as published, not as a live comparator.
  \item \textbf{No bootstrap or subgroup outputs.} We do not report
    bootstrap uncertainty, year trends, or subreddit comparisons.
  \item \textbf{No human raters.} The whole point of our framework is
    that you can validate AI output without humans, but human review
    by clinicians and patients would still be the gold standard.
\end{itemize}

\section{What's next}

The completed run is now ready for manuscript cleanup. The next steps are:

\begin{enumerate}
  \item Pick a target journal (likely JMIR, npj Digital Medicine, or
    a clinical endometriosis journal)
  \item Eventually do a follow-up study where clinicians and patients
    rate the same themes, to validate our automated reliability score
  \item Re-run the analysis later with a second extractor, prompt-B,
    temperature stability, bootstrap uncertainty, and subgroup analyses
    if resources allow
\end{enumerate}

\section{Why a 10th grader should care}

Three reasons:

\begin{enumerate}
  \item \textbf{This is what real-world AI research looks like.} Not
    chatbots. Not deepfakes. People taking massive amounts of public
    text and trying to extract honest insights from it.
  \item \textbf{Endometriosis affects roughly 1 in 10 women}, and
    they've been telling the medical system for decades that they're
    not being heard. AI summaries of their words don't fix that, but
    they at least amplify the signal.
  \item \textbf{Trust-but-verify is the whole game with AI.} Anyone
    can ask an AI a question. Knowing whether the answer is honest is
    the harder problem. The checks we completed are one way to do that.
\end{enumerate}

\bibliographystyle{plain}
\bibliography{refs}

\end{document}
